\chapter{Transport and Coherence of Composition}\label{chap:4}

\section*{Overview of the Chapter}

The question of this chapter is \textbf{under what conditions structure can be carried along
a change of reading and, when such changes are performed in succession, the same result can
be reached coherently}.
Chapter~\ref{chap:1} defined objects and morphisms at the three levels of extraction, core,
and geometry. In this chapter we construct, from a change of extraction and an object before
the change, the core and the geometry after the change.

For example, consider a refactoring that tidies up the field names and data formats of an
order-processing API. We want to read and write the same orders after the change, and to
keep conditions such as ``the total amount equals the sum of the line items.'' Transport in
the sense of this chapter is a construction that carries over to the target not only the
correspondence of data but also the operations and conditions.

Moreover, if there are two routes, one that moves from the old version to the new one through an
intermediate version and one that moves directly, we want the same order to give the
same result along either route. When the conversions are implemented by different teams,
however, discrepancies can arise, such as the correspondence between shipping and billing
addresses being swapped along one route only. Even if each conversion can be undone, the
results of the two routes need not agree.

In the first half we characterize transport by the universal property that every further
change factors through it uniquely. This uniqueness makes the canonical comparisons for
units and composition coherent. In the second half we express the difference between a
separately specified comparison and the canonical comparison as an automorphism, and prove
that this difference can be removed by reselecting edges if and only if all comparison
diagrams can be made commutative at once.

\begin{xltabular}{\linewidth}{@{}LLL@{}}
\toprule
Question & Construction & Outcome \\
\midrule
\endhead
What can be carried from a change of extraction? & Reindexing by a bijection of Atoms & Transport of the core and unique factorization \\
\addlinespace[3pt]
Can covers and coefficients also be carried? & Images of coverage requirements, comparison of overlaps and local data & Transport of the geometry and conditions for its existence \\
\addlinespace[3pt]
What happens under successive transports? & Composition comparisons built from universality & Coherence for units, triple composites, and between levels \\
\addlinespace[3pt]
Can discrepancies of specified comparisons be removed? & Edge reselection and noncommutative defects & Equivalence with the existence of a coherent reselection \\
\addlinespace[3pt]
Can coherence at the levels be combined? & Projection to the core and its kernel & Decision using lifts and the same selection of edges \\
\addlinespace[3pt]
How to preserve the meaning of operations? & Commutative diagrams for lenses and protocols & Simultaneous preservation of reads, updates, and named executions \\
\bottomrule
\end{xltabular}

The diagnosis of Chapter~\ref{chap:3} compared complexes built from Law values and their
classes. The transport of this chapter compares object formation, operations, equations, and
local geometry all at once. What the two have in common is not the statement that equal
values or isomorphic objects exist, but the study of actual comparison maps built from the
input.

\section{Characterizing transport by unique factorization}\label{sec:4.1}

We write the projections of Chapter~\ref{chap:1} as
\begin{equation*}
E_{\mathrm{geom}}\xrightarrow{p}E_{\mathrm{core}}
\xrightarrow{q}B
\tag{4.1}\label{eq:4.1}
\end{equation*}
An object of the base is a pair of an extraction doctrine and a selected source, and a
morphism of the base is an exact change in the sense of Definition~\ref{def:1.27}. The
vocabulary of Atoms is fixed.

\begin{definition}[Strongly opcartesian morphism]\label{def:4.1}

Take a functor $r:E\to B$ and a morphism $\eta:X\to X'$, and let
$r(\eta)=\sigma:b\to b'$.
If, for every object $Y$, every base morphism $\tau:b'\to r(Y)$, and every morphism
$h:X\to Y$ with $r(h)=\tau\sigma$, we have
\begin{equation*}
\exists!\,k:X'\longrightarrow Y,\qquad
r(k)=\tau,\quad k\eta=h
\tag{4.2}\label{eq:4.2}
\end{equation*}
then $\eta$ is called a strongly opcartesian morphism with respect to $r$.
A morphism whose base image is an identity, as in $r(k)=\mathrm{id}_{b'}$, is called a
vertical morphism.

Equation~\eqref{eq:4.2} expresses that, once the base change $\sigma$ has been carried out
first, any further change $h$ splits uniquely, leaving the remaining morphism $k$. Here $\tau$
is arbitrary, which includes the case where the base changes again afterwards. This is the
dual of a strongly cartesian morphism
(\cite[\href{https://stacks.math.columbia.edu/tag/02XJ}{Tag 02XJ, Definition 4.33.1}]{Stacks}).

\end{definition}

\begin{lemma}[Uniqueness and composition of transport]\label{lem:4.2}

If strongly opcartesian morphisms $\eta:X\to X'$ and $\eta':X\to X''$ lift the same base
morphism $\sigma$ from the same object $X$, then there exists a unique isomorphism
\begin{equation*}
c:X'\xrightarrow{\sim}X'',\qquad
r(c)=\mathrm{id},\quad c\eta=\eta'
\tag{4.3}\label{eq:4.3}
\end{equation*}
Moreover, identity morphisms are strongly opcartesian, and composites of such morphisms are
again strongly opcartesian.

\end{lemma}

\begin{proof}

Applying \eqref{eq:4.2} to the identity morphism of the base, we obtain $c$ and a morphism
$d$ in the opposite direction. Applying uniqueness to $dc\eta=\eta$ and $cd\eta'=\eta'$
gives $dc=\mathrm{id}$ and $cd=\mathrm{id}$.
For composites, a morphism out of $X\xrightarrow{\eta}X'\xrightarrow{\zeta}X''$ is factored
first through $\eta$ and then through $\zeta$. The two applications of uniqueness give
uniqueness for the composite. For an identity morphism, $k=h$.

\end{proof}

\section{Constructing the core along an exact change}\label{sec:4.2}

To obtain the transported core, we reread the names of the Atoms as their names after the
change. An input is first sent back to the original names, the original construction is
used, and the output is sent to the new names. The same procedure applies to operations and
to equations.

\begin{construction}[Reindexing of Atoms]\label{cons:4.3}

Let $\sigma=(f,e):(D,s)\to(D',s')$ be an exact morphism of the base. Here $e$ is a
bijection of Atoms and $s'=f(s)$. For families and configurations we use the direct image
$e_*$ of Chapter~\ref{chap:1}. On objects we define
\begin{equation*}
T_e(A)=(e_*C_A,S_A,Q_A)
\tag{4.4}\label{eq:4.4}
\end{equation*}
The structure data and the observed values are kept, and the configuration is carried. The
map $T_{e^{-1}}$ is inverse to $T_e$.
The Atom map $d$ of a morphism of configurations is carried to $ede^{-1}$.

For a core reading $P$, the reading after the change is defined as follows.
\begin{equation*}
\begin{aligned}
\mathrm{Comp}^{\sigma}(F)
 &=e_*\mathrm{Comp}_P(e_*^{-1}F),\\
\mathrm{Form}^{\sigma}(C)
 &=T_e\mathrm{Form}_P(e_*^{-1}C),\\
\mathrm{Op}^{\sigma}(A,B)
 &=\mathrm{Op}_P(T_{e^{-1}}A,T_{e^{-1}}B).
\end{aligned}
\tag{4.5}\label{eq:4.5}
\end{equation*}
The names of the operations are kept as they are, and their actions on configurations are
conjugated. The evaluations of invariants and signatures are defined by composition with
$T_{e^{-1}}$, and the indices, value ranges, and selected axes are kept.

The sets of supports, axes, and observables of a context $W$ are also kept. The condition
that a support reads an Atom $a'$ is defined as the condition that the original support
reads $e^{-1}(a')$.
This yields an equivalence of context categories $\theta_e:\mathcal{C}_P\simeq \mathcal{C}^\sigma$.
We write $\bar\theta_e$ for the reindexing in the opposite direction.
Keeping the equation indices and roles, we define
\begin{equation*}
\begin{aligned}
O^\sigma(W')&=O_P(\bar\theta_eW'),\\
\nu^\sigma_{W',i,a'}&=\nu_{P,\bar\theta_eW',i,e^{-1}(a')},\\
\varepsilon^\sigma_{W',A',i,a'}
 &=\varepsilon_{P,\bar\theta_eW',T_{e^{-1}}A',i,e^{-1}(a')}
\end{aligned}
\tag{4.6}\label{eq:4.6}
\end{equation*}
The restriction maps are also reindexed along $\bar\theta_e$.
In the detector code, we apply $e$ to each Atom appearing in a query.
The positive and negative signs, the logical operations, and the operation names are
unchanged.

\end{construction}

\begin{proposition}[The transported core and the canonical morphism]\label{prop:4.4}

Construction~\ref{cons:4.3} yields a core $\sigma_!P$ over the base $(D',s')$ and a
morphism of cores
$\eta_{\sigma,P}:P\to\sigma_!P$ over the base morphism $\sigma$.
Its extracted family, base object, and generated object family satisfy
\begin{equation*}
F_{\sigma_!P}=e_*F_P,\qquad
A_{\sigma_!P}=T_eA_P,\qquad
\mathrm{Obj}_{\sigma_!P}=T_e(\mathrm{Obj}_P)
\tag{4.7}\label{eq:4.7}
\end{equation*}
\end{proposition}

\begin{proof}

The first equality follows from preservation and reflection of extraction and from the
bijectivity of $e$. Hence the extracted family is finite, and \eqref{eq:4.5} can be
applied. Applying the object-formation formula next gives the equality for the base
objects. Finite sequences of operations are carried into each other by $T_e$ and
$T_{e^{-1}}$, so the equality for the generated object families follows from
Theorem~\ref{thm:1.18}.

The object map of the canonical morphism is $T_e$, and its operation map is the
correspondence of \eqref{eq:4.5}.
The commutative square for configurations is $(ede^{-1})e=ed$.
Equation~\eqref{eq:4.6} preserves coordinates and residuals, and commutation with the
restrictions follows from the original commutation.
Agreement of queries is preserved and reflected for both the positive and the negative
sign, so acceptance and soundness of the detectors are carried as well.
The conditions on invariants and signatures are obtained by composing with the inverse
maps. This verifies all the conditions of Definition~\ref{def:1.28}.

\end{proof}

\begin{theorem}[Universality of core transport]\label{thm:4.5}

The morphism $\eta_{\sigma,P}$ is strongly opcartesian with respect to
$q:E_{\mathrm{core}}\to B$.

\end{theorem}

\begin{proof}

Take $h:P\to R$ with $q(h)=\tau\sigma$.
We build the factor $k$ by sending each datum of $\sigma_!P$ back to its original names
and then applying $h$. For example, its object map is
\begin{equation*}
H_k=H_hT_{e^{-1}}
\tag{4.8}\label{eq:4.8}
\end{equation*}
and the operation map likewise returns to the original operations by the correspondence of
\eqref{eq:4.5} and then applies $\Phi_h$.
The comparison of equations is reindexed along $\bar\theta_e$, and $e^{-1}$ is substituted
into the Atoms of the coordinates, the residuals, and the detectors.
For the maps of invariant indices and of signature axes and value ranges we use those of
$h$. For the base morphism we use $\tau$.

Substituting these into \eqref{eq:4.5}--\eqref{eq:4.6}, the preservation equations for $h$
become exactly the preservation equations for $k$. Hence $k$ is a morphism of cores.
From $T_{e^{-1}}T_e=\mathrm{id}$ and the inverse correspondences of the indices we obtain
$k\eta_{\sigma,P}=h$.

Conversely, the object map of any morphism satisfying this composition equation is
determined to be \eqref{eq:4.8} by the surjectivity of $T_e$.
The components for operations, contexts, coordinates of equations, and invariants and
signatures are likewise determined uniquely through the invertible correspondences of the
canonical morphism.
Since the base is fixed to $\tau$, the whole morphism is unique.

\end{proof}

Bijectivity of the source map $f$ is not needed for this construction.
What is needed is that extraction is preserved and reflected by the specified
correspondence of Atoms.

\begin{example}[Forward preservation of extraction is not enough]\label{ex:4.6}

Let the Atoms be $\{a,b,c\}$, with a single source and identity normalization.
Let the extracted family before the change be $\{a,b\}$ and the one after the change be
$\{a,b,c\}$, and let the correspondence of Atoms be the bijection that swaps $a,b$ and
fixes $c$.
Extraction is preserved in the forward direction, but reflection fails for $c$.
The family after the change is not the direct image of the original family, so this is not
an exact transport satisfying \eqref{eq:4.7}.

Object formation that uses the added Atom requires data about the added part.
The next construction makes this role explicit.

\end{example}

\begin{construction}[Additional data for a change that enlarges extraction]\label{cons:4.7}

Take a change $(f,e)$ that preserves extraction in the forward direction, with $e$ a
bijection.
Finiteness of the extracted family $F'$ after the change is given as part of the data.
Let $A'$ be the base object built from $F'$ by the object formation of \eqref{eq:4.5}, and
add the following data.

\begin{enumerate}
\item An actual operation $A'\to T_eA_P$ in the transported operation reading.
\item An equation system on the contexts of $A'$, together with a bijection to the original
equation indices.
\item Soundness, with respect to that equation system, of the code obtained by carrying the
original detector code along $e$.
\end{enumerate}

From these data the core after the change can be constructed.
Moreover, the images of the original generated objects are contained in the generated
objects after the change.
Indeed, it suffices to follow the added operation by the image of an original finite
sequence of operations.
On the images of the original objects, agreement of queries and acceptance by the
detectors are preserved.
This follows from the commutation of the evaluation of queries and codes with the
reindexing of Atoms.
Operations, invariants, and signatures are also preserved by the corresponding maps.
What this construction provides is a comparison equipped with the forward preservation
just described.

\end{construction}

\section{Lifting the transport of the core to the geometry}\label{sec:4.3}

After the core has been carried, we carry the choice of which local data are viewed
together. For the coverage requirements we use their images, and for the coefficients and
the presentations of rings we use reindexing. The supports, axes, and observables inside a
context are compared together with their readable elements and restriction maps.

\begin{construction}[Transport of the components of a geometry]\label{cons:4.8}

Take a geometry $G$ and a morphism of cores $h:pG\to P'$.
Let $\theta$ be its equivalence of contexts and $\bar\theta$ a specified quasi-inverse.
The coverage requirements after the change are defined as the images under existential
quantification along the maps of Definition~\ref{def:1.30}.
For example, the required supports and their visibility are
\begin{equation*}
\begin{aligned}
\mathrm{Req}'(a')
 &\Longleftrightarrow \exists a,\ \mathrm{Req}(a)\land e(a)=a',\\
\mathrm{Vis}'(W',a')
 &\Longleftrightarrow
 \exists W,a,\ \mathrm{Vis}(W,a)\land\theta W=W'\land e(a)=a'
\end{aligned}
\tag{4.9}\label{eq:4.9}
\end{equation*}
The equation coordinates, the witnesses, the axes, and the visibility between two contexts
are carried by the same rule.
The overlaps are
\begin{equation*}
\mathrm{Ov}'(U'\to W'\leftarrow V')
=\theta\bigl(\mathrm{Ov}(\bar\theta U'\to
\bar\theta W'\leftarrow\bar\theta V')\bigr)
\tag{4.10}\label{eq:4.10}
\end{equation*}
and the projections are sent to $U',V'$ using the counit.
The universal property of the overlaps follows from the equivalence $\theta,\bar\theta$.
The coefficient ring $k$ is kept, and the raw system is
$\mathcal{B}'=\bar\theta^*\mathcal{B}$.
The coordinate types, the local data types, and the relation polynomials and restrictions
are all reindexed.
We write $h_*G$ for this geometry.

\end{construction}

\begin{definition}[Transport condition for local realizations]\label{def:4.9}

Define $H_{\mathrm{geom}}(G,h)$ to be a family, over the contexts, of three comparison maps
\begin{equation*}
\begin{aligned}
s_W&:\mathrm{Supp}(W)\to\mathrm{Supp}'(\theta W),\\
a_W&:\mathrm{Ax}(W)\to\mathrm{Ax}'(\theta W),\\
o_W&:\mathrm{Obs}(W)\to\mathrm{Obs}'(\theta W)
\end{aligned}
\tag{4.11}\label{eq:4.11}
\end{equation*}
that preserves readable elements in the sense of Definition~\ref{def:1.30} and satisfies
the naturality of \eqref{eq:1.14}.
The three sets on the target side and their readouts are those carried by the contexts of
the target core.

\end{definition}

\begin{proposition}[Existence condition for morphisms of geometries]\label{prop:4.10}

The data of $H_{\mathrm{geom}}(G,h)$ exist if and only if a morphism of geometries
$G\to h_*G$ exists over $h$.
Moreover, any morphism of geometries over $h$ yields $H_{\mathrm{geom}}(G,h)$.

\end{proposition}

\begin{proof}

Equation~\eqref{eq:4.9} gives forward preservation of the coverage requirements.
Equation~\eqref{eq:4.10} gives the comparison of overlaps; for the coefficients we use the
identity map, and for the raw system the reindexing equality.
The remaining comparisons are \eqref{eq:4.11} together with their preservation and
naturality, and these are exactly the data of $H_{\mathrm{geom}}(G,h)$.
In the opposite direction, one extracts this family of three comparisons from a morphism
of geometries.

\end{proof}

\begin{theorem}[Canonical geometry transport and universality]\label{thm:4.11}

For the canonical morphism of cores $h=\eta_{\sigma,pG}$, the data
$H_{\mathrm{geom}}(G,h)$ can be constructed.
We write
$\sigma_!^{\mathrm g}G$ and $\widetilde\eta_{\sigma,G}:G\to\sigma_!^{\mathrm g}G$
for the resulting geometry and morphism. This morphism is strongly opcartesian with
respect to $p:E_{\mathrm{geom}}\to E_{\mathrm{core}}$, and it satisfies
\begin{equation*}
p(\sigma_!^{\mathrm g}G)=\sigma_!(pG),\qquad
p(\widetilde\eta_{\sigma,G})=\eta_{\sigma,pG}
\tag{4.12}\label{eq:4.12}
\end{equation*}
\end{theorem}

\begin{proof}

In the context transport of Construction~\ref{cons:4.3}, the sets of supports, axes, and
observables were kept. We use these identity correspondences in \eqref{eq:4.11}. Since the
readout of supports was defined by $e^{-1}$, the condition of mapping Atoms by $e$ is
satisfied, and the three naturalities follow from the original structure maps.
Hence Proposition~\ref{prop:4.10} gives a morphism of geometries.

To prove universality, take an arbitrary morphism of cores
$t:\sigma_!(pG)\to pK$ and a morphism of geometries $F:G\to K$ with
$p(F)=t\eta_{\sigma,pG}$.
We send each context back by $\bar\theta_e$, invert the three invertible correspondences
of the canonical morphism, and then apply the comparison maps of $F$.
This defines the comparisons of supports, axes, and observables of the factor.
The coefficient map is set equal to the coefficient map of $F$.

Each coverage requirement on the target has, by \eqref{eq:4.9}, a witness coming from an
original requirement.
Applying the preservation laws of $F$ to that witness yields the preservation laws of the
factor.
The comparison of overlaps is built by composing the two equivalences of categories, and
the equality of raw systems is obtained from the commutation of reindexing with the change
of coefficients.
We thus obtain a morphism of geometries whose base is $t$, and its composite with the
canonical morphism is $F$.

The three comparisons of the canonical morphism are bijections and its coefficient map is
the identity, so each map of the factor is unique.
The comparisons of overlaps inside the context category are determined by the uniqueness
of parallel morphisms in a preorder category.
This establishes the existence and uniqueness of \eqref{eq:4.2}.
Equation~\eqref{eq:4.12} follows from the construction.

\end{proof}

For a general morphism of cores, Proposition~\ref{prop:4.10} gives the existence
condition; for the canonical core transport, both that condition and the universal
property have now been constructed.

\section{Units, composition, and projections of transport}\label{sec:4.4}

The universal property of transport determines not only how objects are carried but also
how comparison morphisms are carried. Carrying morphisms between objects over the same
base yields functors between the fibers.

\begin{lemma}[Universality in the tower]\label{lem:4.12}

If a morphism $v:G\to G'$ is strongly opcartesian with respect to $p$, and $p(v)$ is
strongly opcartesian with respect to $q$, then
$v$ is strongly opcartesian with respect to $qp$.

\end{lemma}

\begin{proof}

A factorization problem over $qp$ is first split by the universal property of $q$ into a
factor at the level of cores, and the universal property of $p$ is then applied with that
core factor as the base morphism.
Existence and uniqueness at the two levels give the required existence and uniqueness.

\end{proof}

\begin{construction}[Transport functors between fibers]\label{cons:4.13}

For a functor $r:E\to B$, the fiber $E_b$ is the category consisting of the
objects satisfying $r(X)=b$ and the vertical morphisms between them.
For each $\sigma:b\to b'$ and $X\in E_b$, choose a strongly opcartesian lift
$\eta_{\sigma,X}:X\to T_\sigma X$.
The image of a vertical morphism $v:X\to Y$ is defined to be the unique vertical morphism
\begin{equation*}
T_\sigma(v):T_\sigma X\longrightarrow T_\sigma Y,\qquad
T_\sigma(v)\eta_{\sigma,X}=\eta_{\sigma,Y}v
\tag{4.13}\label{eq:4.13}
\end{equation*}
The right-hand side lies over the base $\sigma$, so this morphism exists.
Identity morphisms and composites of two morphisms satisfy the same equation, so
uniqueness shows that $T_\sigma:E_b\to E_{b'}$ is a functor.
By Theorems~\ref{thm:4.5} and~\ref{thm:4.11} and Lemma~\ref{lem:4.12}, this
construction applies to cores and to geometries with base $B$.

\end{construction}

\begin{construction}[Composition comparisons and unit comparisons]\label{cons:4.14}

For $b\xrightarrow{\sigma}b'\xrightarrow{\tau}b''$, applying Lemma~\ref{lem:4.2} to the
twofold lift and the single lift yields an isomorphism
\begin{equation*}
\begin{aligned}
c_{\tau,\sigma,X}&:T_\tau T_\sigma X\xrightarrow{\sim}T_{\tau\sigma}X,\\
c_{\tau,\sigma,X}\eta_{\tau,T_\sigma X}\eta_{\sigma,X}
 &=\eta_{\tau\sigma,X}
\end{aligned}
\tag{4.14}\label{eq:4.14}
\end{equation*}
For identity morphisms we similarly obtain
\begin{equation*}
\epsilon_X:T_{\mathrm{id}}X\xrightarrow{\sim}X,\qquad
\epsilon_X\eta_{\mathrm{id},X}=\mathrm{id}_X
\tag{4.15}\label{eq:4.15}
\end{equation*}
We call these the compositor and the unitor, respectively.

\end{construction}

\begin{theorem}[Coherence of composition comparisons]\label{thm:4.15}

The families $c_{\tau,\sigma}$ and $\epsilon$ are natural isomorphisms.
Moreover, for composable morphisms $\sigma,\tau,\rho$,
\begin{equation*}
\begin{aligned}
c_{\rho,\tau\sigma,X}\,T_\rho(c_{\tau,\sigma,X})
 &=c_{\rho\tau,\sigma,X}\,c_{\rho,\tau,T_\sigma X},\\
c_{\mathrm{id},\sigma,X}&=\epsilon_{T_\sigma X},\\
c_{\sigma,\mathrm{id},X}&=T_\sigma(\epsilon_X)
\end{aligned}
\tag{4.16}\label{eq:4.16}
\end{equation*}
Associativity and units of composed base morphisms are identified through the equalities
of the category $B$.
The structure consisting of the fibers, the transport functors, and these natural
isomorphisms is called a pseudofunctor.

\end{theorem}

\begin{proof}

We compose the two routes of the naturality square for $c$ with the twofold lift from the
starting point.
By \eqref{eq:4.13}--\eqref{eq:4.14}, both become $\eta_{\tau\sigma,Y}v$.
The twofold lift is strongly opcartesian, so the two routes are equal.
Naturality of $\epsilon$ follows by the same argument.

Both sides of the first equation of \eqref{eq:4.16} are vertical morphisms from
$T_\rho T_\tau T_\sigma X$ to $T_{\rho\tau\sigma}X$.
Composed with the threefold lift, both become $\eta_{\rho\tau\sigma,X}$.
Again they are equal by uniqueness. The two unit equations are checked in the same way by
composing with the corresponding twofold lifts.

\end{proof}

This canonical construction is the dual of the cartesian case treated in
\cite[\href{https://stacks.math.columbia.edu/tag/02XJ}{Tag 02XJ, Lemmas 4.33.2 and 4.33.7}]{Stacks}.
The reason for the coherence lies in the uniqueness of the factorizations that determined
the comparisons.

\begin{theorem}[Compatibility with the projection to the core]\label{thm:4.16}

Let $p_b$ be the projection from the fiber of geometries to the fiber of cores.
Write $T^{\mathrm g}_\sigma$ and $T^{\mathrm c}_\sigma$ for the transport of geometries
and the transport of cores, respectively.
There is a natural isomorphism
\begin{equation*}
\chi_\sigma:p_{b'}T^{\mathrm g}_\sigma
\xRightarrow{\sim}T^{\mathrm c}_\sigma p_b
\tag{4.17}\label{eq:4.17}
\end{equation*}
and, for composition and units, it satisfies
\begin{equation*}
\begin{aligned}
\chi_{\tau\sigma,G}\,p(c^{\mathrm g}_{\tau,\sigma,G})
 &=
c^{\mathrm c}_{\tau,\sigma,pG}\,
T^{\mathrm c}_\tau(\chi_{\sigma,G})\,
\chi_{\tau,T^{\mathrm g}_\sigma G},\\
p(\epsilon^{\mathrm g}_G)
 &=\epsilon^{\mathrm c}_{pG}\,\chi_{\mathrm{id},G}
\end{aligned}
\tag{4.18}\label{eq:4.18}
\end{equation*}
\end{theorem}

\begin{proof}

The morphisms $p(\widetilde\eta_{\sigma,G})$ and $\eta_{\sigma,pG}$ are strongly
opcartesian morphisms lifting the same base morphism from the same object.
The unique comparison connecting them is $\chi_{\sigma,G}$.
In the canonical construction the two lifts correspond by \eqref{eq:4.12}.
Equation~\eqref{eq:4.13} still holds after projecting a vertical morphism of geometries,
so the comparison is natural.
The two routes of \eqref{eq:4.18} also become the same morphism when composed with the
lift from the original core.
Uniqueness gives the two equations.

\end{proof}

Hence the route that carries the geometry and then reads off the core and the route that
reads off the core and then carries it are compatible, including objects, morphisms, and
composition comparisons.

\section{The difference between specified and canonical comparisons}\label{sec:4.5}

The canonical comparisons were coherent by uniqueness.
In practice, however, a correspondence of state names or of local data is sometimes
specified separately.
To measure how far such a correspondence differs from the canonical comparison, we read
both as automorphisms of the same object.

Henceforth, let $r:E\to B$ be a functor and choose strongly opcartesian
morphisms on the required edges.
For cores we may use $r=q$, and for full geometries, $r=qp$.

\begin{definition}[Finite comparison data]\label{def:4.17}

On each vertex $i$ of a finite directed graph we place an object $X_i$, and on each edge
$e:i\to j$ a base morphism $\sigma_e:r(X_i)\to r(X_j)$ and a strongly opcartesian
morphism $L_e:X_i\to X_j$ over it.
We write $L_\pi$ for the composite along a path $\pi$; the composite along the empty path
is the identity.

We further specify finitely many faces $f$.
A face is a pair of paths $\ell_f,\rho_f:i\to j$ with the same start and end whose
composites in the base are equal.
To each face we assign a fiber automorphism of the end,
$u_f\in\mathrm{Aut}_r(X_j)$, where
\begin{equation*}
\mathrm{Aut}_r(X)
=\{a:X\xrightarrow{\sim}X\mid r(a)=\mathrm{id}_{r(X)}\}
\tag{4.19}\label{eq:4.19}
\end{equation*}
The automorphism $u_f$ represents the comparison we wish to adopt on that face.
At this stage we do not assume $u_fL_{\ell_f}=L_{\rho_f}$.

\end{definition}

\begin{construction}[Canonical comparisons and raw defects]\label{cons:4.18}

By Lemma~\ref{lem:4.2}, each face carries a unique fiber automorphism
\begin{equation*}
\phi_fL_{\ell_f}=L_{\rho_f},\qquad
\phi_f\in\mathrm{Aut}_r(X_j)
\tag{4.20}\label{eq:4.20}
\end{equation*}
The difference from the specified comparison is defined as
\begin{equation*}
\delta_f=u_f\phi_f^{-1}
\tag{4.21}\label{eq:4.21}
\end{equation*}
and called the raw defect.
Thus $\delta_f=1$ means that the specified comparison agrees with the canonical one.
Products are composites of morphisms, and the right-hand factor acts first.

If the automorphism group is noncommutative, the result depends on the order of the
comparisons.
For this reason a discrepancy cannot in general be treated as a numerical difference or as
an element of a commutative group.
Equation~\eqref{eq:4.21} is a difference that also records the order of the comparisons.

\end{construction}

\section{Removing discrepancies by reselecting edges}\label{sec:4.6}

We now keep the comparisons fixed and change the lift adopted on each edge.
By Lemma~\ref{lem:4.2}, lifts of the same base morphism with the same start and end are
carried into one another by composing with a fiber automorphism of the end.

\begin{definition}[Edge reselection]\label{def:4.19}

For each edge $e:i\to j$ choose
$a_e\in\mathrm{Aut}_r(X_j)$ and set
\begin{equation*}
L_e^a=a_eL_e,\qquad
\mathcal{A}=\prod_{e:i\to j}\mathrm{Aut}_r(X_j)
\tag{4.22}\label{eq:4.22}
\end{equation*}
We call $a=(a_e)_e$ an edge reselection.
A vertical isomorphism $a_e$ is strongly opcartesian, since the factor in \eqref{eq:4.2}
is uniquely determined as $k=ha_e^{-1}$. By Lemma~\ref{lem:4.2}, the composite $a_eL_e$ is
then strongly opcartesian as well.
For the composites $L_\pi^a$ along paths we define \eqref{eq:4.20}--\eqref{eq:4.21} in the
same way and write $\phi_f(a)$ and $\delta_f(a)$.
The specified comparisons $u_f$ stay fixed.

\end{definition}

\begin{example}[A comparison of two independent edges can be absorbed into one of them]\label{ex:4.20}

Regard a group as a one-object category, with the terminal category as base.
All morphisms are invertible and strongly opcartesian.
For the group product we use composition acting from the right.
Take two parallel edges $x,y$ and one face $x\Rightarrow y$.
Let both original lifts be the identity, and let the specified comparison be an arbitrary
group element $u$.
For edge selections $a_x,a_y$ we have
\begin{equation*}
\phi(a)=a_ya_x^{-1},\qquad
\delta(a)=u\,a_xa_y^{-1}
\tag{4.23}\label{eq:4.23}
\end{equation*}
Setting $a_x=1$ and $a_y=u$ makes the defect vanish.
The two edges can be reselected independently, so the specified comparison can be absorbed
into one of the edges.

In a general comparison diagram, the selection at each
edge acts on several paths and faces.
We therefore examine the change arising at the end of a path, and determine how it alters
the comparison on each face.

\end{example}

\begin{lemma}[The change at the end of a path]\label{lem:4.21}

If a reselection $a$ is followed by $b$, each path $\pi:i\to j$ carries a unique
automorphism
\begin{equation*}
A_\pi(a,b)L_\pi^a=L_\pi^{ba}
\tag{4.24}\label{eq:4.24}
\end{equation*}
where $(ba)_e=b_ea_e$.
This automorphism satisfies
\begin{equation*}
\begin{aligned}
A_\pi(a,1)&=1,\\
A_\pi(ba,c)A_\pi(a,b)&=A_\pi(a,cb)
\end{aligned}
\tag{4.25}\label{eq:4.25}
\end{equation*}
and, for each face,
\begin{equation*}
\begin{aligned}
\phi_f(ba)
 &=A_{\rho_f}(a,b)\phi_f(a)A_{\ell_f}(a,b)^{-1},\\
\delta_f(ba)
 &=\bigl(u_fA_{\ell_f}(a,b)u_f^{-1}\bigr)
   \delta_f(a)A_{\rho_f}(a,b)^{-1}
\end{aligned}
\tag{4.26}\label{eq:4.26}
\end{equation*}
\end{lemma}

\begin{proof}

Existence and uniqueness in \eqref{eq:4.24} follow from the composite along a path being
strongly opcartesian.
Composing both sides of \eqref{eq:4.25} with $L_\pi^a$ gives the same morphism, so
uniqueness applies again.
The first equation of \eqref{eq:4.26} is obtained by applying uniqueness to
$\phi_f(ba)A_{\ell_f}(a,b)L_{\ell_f}^a=A_{\rho_f}(a,b)\phi_f(a)L_{\ell_f}^a$.
Substituting its inverse into \eqref{eq:4.21} yields the second.

\end{proof}

\begin{proposition}[The action of reselections]\label{prop:4.22}

Let $\mathcal{Z}=\prod_f\mathrm{Aut}_r(X_{\mathrm{end}(f)})$ be the family of
automorphisms at the ends of the faces.
On $\mathcal{A}\times \mathcal{Z}$, a left action is defined by
\begin{equation*}
\Gamma_b(a,z)=
\left(ba,\
\left(
u_fA_{\ell_f}(a,b)u_f^{-1}\,
z_f\,A_{\rho_f}(a,b)^{-1}
\right)_f\right)
\tag{4.27}\label{eq:4.27}
\end{equation*}
In particular, the orbit of $(1,\delta(1))$ is
$\{(a,\delta(a))\mid a\in \mathcal{A}\}$.

\end{proposition}

\begin{proof}

Substituting \eqref{eq:4.25} gives $\Gamma_1=\mathrm{id}$ and
$\Gamma_c\Gamma_b=\Gamma_{cb}$.
The description of the orbit is by \eqref{eq:4.26}.

\end{proof}

This action also records the current edge selection $a$.
The reason is that, for a path through several edges, the change $A_\pi(a,b)$ at the end
depends on the current selection.

For example, take the functor from a one-object group category to the terminal
category, and let both initial edge lifts be identities.
Writing $12$ for the path that traverses edge 1 and then edge 2, we have
\[
\begin{aligned}
L_{12}^{a}&=a_2a_1,\\
L_{12}^{ba}&=b_2a_2b_1a_1,\\
A_{12}(a,b)&=L_{12}^{ba}(L_{12}^{a})^{-1}
=b_2a_2b_1a_2^{-1}.
\end{aligned}
\]
The current selection $a_2$ remains in the last expression.
In a noncommutative group, even reselection by the same $b_1,b_2$ can produce
different changes at the end depending on $a_2$.

\begin{theorem}[Vanishing of the obstruction and coherent selections]\label{thm:4.23}

We call the orbit of Proposition~\ref{prop:4.22} the obstruction class of the finite
comparison data.
When the orbit meets $\mathcal{A}\times\{1\}$, we say that the obstruction vanishes.
This vanishing is equivalent to the existence of an edge reselection $a$ such that, for
all faces simultaneously,
\begin{equation*}
u_fL_{\ell_f}^a=L_{\rho_f}^a
\tag{4.28}\label{eq:4.28}
\end{equation*}
\end{theorem}

\begin{proof}

Vanishing of the obstruction amounts to $\delta_f(a)=1$ for all faces, which by
\eqref{eq:4.21} is equivalent to $u_f=\phi_f(a)$.
This equality gives \eqref{eq:4.28}.
Conversely, if \eqref{eq:4.28} holds, then $u_f$ and $\phi_f(a)$ both send
$L_{\ell_f}^a$ to $L_{\rho_f}^a$.
They are equal by the uniqueness for strongly opcartesian morphisms.

\end{proof}

\section{Additional conditions for pasting faces}\label{sec:4.7}

Performing the comparison of one face and then that of another produces a route that
rewrites paths.
For example, when three paths are compared, there are the direct comparison and the
comparison through the third path.
Agreement of these two routes must also be examined as a separate condition.

\begin{lemma}[Transport of automorphisms along a subsequent path]\label{lem:4.24}

Fix an edge selection $a$.
For a path $\beta:j\to k$ and $v\in\mathrm{Aut}_r(X_j)$, a unique automorphism
$\beta_*v$ satisfies
\begin{equation*}
(\beta_*v)L_\beta^a=L_\beta^av
\tag{4.29}\label{eq:4.29}
\end{equation*}
The map $\beta_*$ is a group homomorphism.
\end{lemma}

\begin{proof}

We use the universal property of the strongly opcartesian morphism $L_\beta^a$.
The morphism obtained from $v^{-1}$ is the inverse.
For identities and products, apply \eqref{eq:4.29} and uniqueness as well.

\end{proof}

\begin{construction}[Pasting of comparisons]\label{cons:4.25}

Consider a single step that rewrites the path $\alpha\ell_f\beta$ into
$\alpha \rho_f\beta$.
Paths proceed from left to right. Since composition of morphisms acts from the right,
\[
L_{\alpha\ell_f\beta}^{a}
=L_\beta^{a}L_{\ell_f}^{a}L_\alpha^{a}
\]
The specified comparison of this step is $\beta_*u_f$, and its canonical comparison is
$\beta_*\phi_f(a)$.
For a rewriting in the opposite direction we use the respective inverses.

Along a sequence of rewritings $P$, we write $U_P(a)$ and $\Phi_P(a)$ for the composites
of these in order.
If the initial path is $\pi$ and the final path is $\rho$, then
\begin{equation*}
\Phi_P(a)L_\pi^a=L_\rho^a,\qquad
D_P(a)=U_P(a)\Phi_P(a)^{-1}
\tag{4.30}\label{eq:4.30}
\end{equation*}
We call $D_P(a)$ the pasted raw defect.
The canonical comparison $\Phi_P(a)$ is uniquely determined by the initial and final
paths.

\end{construction}

For two successive rewrites, write $u_1,u_2$ for the specified comparisons,
$\phi_1,\phi_2$ for the canonical comparisons, and $\delta_1,\delta_2$ for the defects.
Then
\begin{equation*}
\begin{aligned}
U_P\Phi_P^{-1}
 &=u_2u_1(\phi_2\phi_1)^{-1}\\
 &=(u_2\delta_1u_2^{-1})\delta_2.
\end{aligned}
\tag{4.31}\label{eq:4.31}
\end{equation*}

The later comparison conjugates the earlier defect.
Thus the defect of a pasting cannot be replaced by the simple product of the defects
of its faces.

\begin{proposition}[Coherence condition among relations]\label{prop:4.26}

A pair of rewriting sequences $P,Q$ from the same path to the same path is called a
syzygy.
Specify finitely many syzygies, and suppose that for the edge selection $a$
\begin{equation*}
U_P(a)=U_Q(a)
\tag{4.32}\label{eq:4.32}
\end{equation*}
holds for each of them. Then
$D_P(a)=D_Q(a)$.

\end{proposition}

\begin{proof}

The comparisons $\Phi_P(a)$ and $\Phi_Q(a)$ send the same strongly opcartesian morphism
to the same morphism, hence are equal. Multiply \eqref{eq:4.32} by their common inverse
on the right.

\end{proof}

A raw defect is defined by \eqref{eq:4.21} alone.
To obtain the cocycle condition $D_P(a)=D_Q(a)$ along the specified syzygies,
\eqref{eq:4.32} is required.
Condition~\eqref{eq:4.32} is a relation among the specified comparisons; it does not
assume the vanishing of the defect on each face.

\begin{proposition}[Obstruction for closed comparisons]\label{prop:4.27}

For any syzygy, writing $\Phi$ for the common canonical comparison,
\begin{equation*}
D_Q^{-1}D_P
=\Phi\,(U_Q^{-1}U_P)\,\Phi^{-1}
\tag{4.33}\label{eq:4.33}
\end{equation*}
In particular, the left-hand side being the identity is equivalent to \eqref{eq:4.32}.

\end{proposition}

\begin{proof}

Substitute $D_P=U_P\Phi^{-1}$ and $D_Q=U_Q\Phi^{-1}$.

\end{proof}

Equation~\eqref{eq:4.33} shows that, once the edge selection is fixed, a mismatch of the
specified pastings persists through conjugation.
If every face satisfies \eqref{eq:4.28}, then every specified pasting agrees with the
canonical comparison.
Failure of \eqref{eq:4.32} therefore obstructs overall coherence for that edge selection.

\section{Removable and irremovable discrepancies}\label{sec:4.8}

In Example~\ref{ex:4.20}, the comparison of a single face could be made coherent by
reselecting one of two independent edges.
When several faces share edges, each face may be made coherent on its own, and yet not all of them at
once by the same selection.
Again, regard a group as a one-object category, with the terminal category as base.

\begin{example}[Different requirements on the same two edges]\label{ex:4.28}

Let the group be the permutation group $S_3$ on three points, and place two faces on the
same two edges.
Let the specified comparisons be $s=(12)$ and $t=(23)$.
The canonical comparison of both faces is the same $a_ya_x^{-1}$, so coherence requires
\begin{equation*}
a_ya_x^{-1}=s,\qquad a_ya_x^{-1}=t
\tag{4.34}\label{eq:4.34}
\end{equation*}
Since $s\ne t$, no selection satisfies both at once.
Each face on its own can be made coherent by Example~\ref{ex:4.20}.
Making the two faces coherent by the same edge selection is the new condition.

\end{example}

\begin{example}[Three comparisons and order]\label{ex:4.29}

Take three parallel edges $x_0,x_1,x_2$, with all original lifts the identity.
Specify the comparison of $x_0\Rightarrow x_1$ to be $s$,
that of $x_1\Rightarrow x_2$ to be $t$, and
that of $x_0\Rightarrow x_2$ to be $st$.
The indirect comparison is $ts$ and the direct comparison is $st$.
For $s=(12)$ and $t=(23)$, the two differ.

Indeed, if the first two faces are coherent, then
$a_1=sa_0$ and $a_2=ta_1=tsa_0$.
The third face requires $a_2=sta_0$, so the three faces are not coherent simultaneously.
The corresponding syzygy has no subsequent path, and the two specified routes do not
depend on the edge selection.
The inner discrepancy in \eqref{eq:4.33} is
\begin{equation*}
(st)^{-1}ts=t^{-1}s^{-1}ts\ne1
\tag{4.35}\label{eq:4.35}
\end{equation*}
and it persists for every selection.
In this example, besides each correspondence being a bijection, the orders of composition
must be brought into line.

\end{example}

\section{Achieving coherence of the core and the geometry with the same selection}\label{sec:4.9}

Projecting a comparison of geometries to the core yields a comparison of cores.
After correcting the discrepancies visible at the core, we want to correct the
discrepancies that remain only in the geometry.
The key that connects these two stages is to lift the changes selected at the core to the
geometry, and to examine the remaining conditions using that same selection.

\begin{definition}[The projection and its kernel]\label{def:4.30}

For an object $G$ and $P=pG$, set
\begin{equation*}
\begin{aligned}
C_G&=\mathrm{Aut}_{qp}(G),&
B_P&=\mathrm{Aut}_q(P),\\
\pi_G:C_G&\longrightarrow B_P,&
\pi_G(a)&=p(a),\\
H_G&=\ker\pi_G=\mathrm{Aut}_p(G)
\end{aligned}
\tag{4.36}\label{eq:4.36}
\end{equation*}
Here $C_G$ represents the changes of the geometry that fix the extraction, $B_P$ the
changes of the core that fix the extraction, and $H_G$ the changes of the geometry that
fix the entire core.
The map $\pi_G$ is a group homomorphism, and $H_G$ is a normal subgroup.

Place $G_i$ on each vertex of Definition~\ref{def:4.17}, and assume that the
edges $L_e$ of geometries and their projections $p(L_e)$ are strongly opcartesian
with respect to $qp$ and $q$, respectively.
The transport constructed in \S\ref{sec:4.4} satisfies this condition.
The specified comparisons are $u_f\in C_{G_j}$, and for the specified comparisons of
cores we use $p(u_f)$.
We write $\phi_f$ for the canonical comparison of geometries and $\bar\phi_f$ for the
canonical comparison of cores.

\end{definition}

\begin{proposition}[Projection of defects]\label{prop:4.31}

For every edge reselection $c$ of the geometries,
\begin{equation*}
p(\phi_f(c))=\bar\phi_f(p(c)),\qquad
p(\delta_f(c))=\bar\delta_f(p(c))
\tag{4.37}\label{eq:4.37}
\end{equation*}
\end{proposition}

\begin{proof}

Project $\phi_f(c)L_{\ell_f}^c=L_{\rho_f}^c$.
The paths of cores are also strongly opcartesian, so the first equation follows from the
uniqueness of comparisons.
The second follows because $p$ preserves composites and inverses.

\end{proof}

\begin{definition}[Edgewise lifts and alignment]\label{def:4.32}

For an edge reselection $a_e\in B_{P_j}$ of the cores, choose
$\widetilde a_e\in C_{G_j}$ with $p(\widetilde a_e)=a_e$.
These finitely many edgewise lifts are called an edge section.
The section is said to be aligned when, for each face, the composites of the core paths
satisfy
\begin{equation*}
p(u_f)\,\bar L_{\ell_f}^{\,a}
=\bar L_{\rho_f}^{\,a}
\tag{4.38}\label{eq:4.38}
\end{equation*}
Here $\bar L^{\,a}$ is the path whose projected edges have been reselected by $a$.

Alignment is the condition that all comparisons visible at the core are coherent for that
edge selection.
As in Definition~\ref{def:4.32}, an edge section is the datum that provides, together
with the reselection at the core, a lift realizing it on the side of the geometry.

\end{definition}

\begin{theorem}[Order-preserving decomposition of the defect]\label{thm:4.33}

Take an aligned edge section, and write $\phi_f^0=\phi_f(1)$ for the comparison before
reselection and $m_f=\phi_f(\widetilde a)$ for the comparison at the chosen section. Then
\begin{equation*}
p(m_f)=p(u_f)
\tag{4.39}\label{eq:4.39}
\end{equation*}
and the original defect decomposes as
\begin{equation*}
\underbrace{u_f(\phi_f^0)^{-1}}_{\delta_f}
=
\underbrace{u_fm_f^{-1}}_{\delta_f^{\mathrm{in}}}
\underbrace{m_f(\phi_f^0)^{-1}}_{\widetilde\delta_f^{\mathrm{out}}}
\tag{4.40}\label{eq:4.40}
\end{equation*}
where
$\delta_f^{\mathrm{in}}\in H_{G_j}$ and
$p(\widetilde\delta_f^{\mathrm{out}})=\bar\delta_f(1)$.

\end{theorem}

\begin{proof}

By Proposition~\ref{prop:4.31} we have $p(m_f)=\bar\phi_f(a)$.
From \eqref{eq:4.38} and the uniqueness of canonical comparisons we obtain
$\bar\phi_f(a)=p(u_f)$, which is \eqref{eq:4.39}.
Hence $p(u_fm_f^{-1})=1$.
Equation~\eqref{eq:4.40} is checked by cancelling the middle $m_f^{-1}m_f$.
The last claim follows from
$p(m_f(\phi_f^0)^{-1})=p(u_f)\bar\phi_f(1)^{-1}$.

\end{proof}

In this decomposition, the factor that lifts the defect visible at the core acts first,
and the factor invisible to the core acts next.
In general the order of the two factors cannot be interchanged.

\begin{theorem}[Condition for simultaneous vanishing]\label{thm:4.34}

Fix an aligned edge section $\widetilde a$.
Then there exists, over its core selection $a$, an edge selection of the geometries that
is coherent on all faces if and only if there exists a family of kernel elements
$h_e\in H_{G_j}$ with
\begin{equation*}
u_f=\phi_f(h\widetilde a)\qquad\text{(for every face }f\text{)}
\tag{4.41}\label{eq:4.41}
\end{equation*}
Consequently, the obstruction for the geometries vanishes if and only if an aligned edge
section can be chosen for which \eqref{eq:4.41} can be satisfied.

\end{theorem}

\begin{proof}

Since $h$ lies in the kernel, $p(h\widetilde a)=a$.
By \eqref{eq:4.41} and Theorem~\ref{thm:4.23}, $h\widetilde a$ is the
required coherent selection.
Conversely, given a coherent selection $c$ with $p(c)=a$, set
$h_e=c_e\widetilde a_e^{-1}$ on each edge.
Its projection is $a_ea_e^{-1}=1$, so $h_e\in H_{G_j}$, and
$c=h\widetilde a$ gives \eqref{eq:4.41}.
For the forward direction of the final equivalence, the coherent selection of the
geometries itself can be used as the edge section.

\end{proof}

Here the same edge selection is used for the core and the geometry, and all faces are
examined at once.
Even if a coherent selection is found at each level on its own, the question remains whether
that selection can be lifted to the next level.

\begin{example}[Four axes and a solution that does not lift]\label{ex:4.35}

We illustrate the distinction by a finite tower in which groups are regarded
as one-object categories.
Regard the four indices $1,2,3,4$ as axes, and take their permutation group $S_4$
as the group of core changes. At the level of geometries, select $\{1,2\}$
among the four axes. From the permutations preserving this selection and a
two-valued extra component, form
\begin{equation*}
\begin{aligned}
K&=\{\gamma\in S_4\mid\gamma\{1,2\}=\{1,2\}\}\times C_2,\\
\pi:K&\longrightarrow S_4,\qquad(\gamma,z)\longmapsto\gamma
\end{aligned}
\tag{4.42}\label{eq:4.42}
\end{equation*}
The group $C_2=\mathbb{Z}/2\mathbb{Z}$ is written additively.
The kernel is $\{1\}\times C_2$, a degree of freedom that changes only the geometry.
Take the terminal category as the base category.

Take one vertex and two loops $e,g$, with the original lifts the identity.
The faces and the specified comparisons are
\begin{equation*}
e^2\Rightarrow\varnothing:\quad u=(b,0),\quad b=(12)(34),
\qquad
g\Rightarrow\varnothing:\quad v=(1,1)
\tag{4.43}\label{eq:4.43}
\end{equation*}
Here $\varnothing$ is the empty path.
The coherence conditions at the core are $a_e^2=b$ and $a_g=1$, and
$a_e=(1324)$, $a_g=1$ is a solution.

On the other hand, the permutations allowed at the geometry preserve $\{1,2\}$ and $\{3,4\}$
separately.
The permutation part is therefore $S_2\times S_2$, in which every element squares to the
identity.
No selection $c_e$ of the geometry can satisfy $c_e^2=(b,0)$.
The square root $(1324)$ at the core does not preserve the two selected axes.

The second face, belonging only to the geometry, is coherent with the kernel selection $c_g=(1,1)$.
Thus all core conditions and the conditions belonging only to the geometry each
have solutions.
Still, no overall solution exists.
The obstruction is the absence of an edge section lifting the core solution.

\end{example}

\section{Carrying the reads and updates of a lens by the same map}\label{sec:4.10}

To connect the comparisons of the preceding sections with software operations, we must
determine which changes preserve the meaning of the operations.
In the lenses of Chapter~\ref{chap:1}, one state appeared both in reads and in updates.
Through this sharing, the maps of the two commutative diagrams are tied together.

\begin{definition}[Morphisms of lenses including a change of view]\label{def:4.36}

A relative morphism between lenses $L=(C,V,g,p)$ and $L'=(C',V',g',p')$
satisfying the three laws is a pair of maps $h:C\to C'$ and $u:V\to V'$ such that
\begin{equation*}
g'h=ug,\qquad
h\,p=p'(h\times u)
\tag{4.44}\label{eq:4.44}
\end{equation*}

Identities and composition are componentwise and preserve \eqref{eq:4.44}.
Taking $u=\mathrm{id}_V$ recovers the fixed-view morphisms of Chapter~\ref{chap:1}.

\end{definition}

Under the lens laws of this chapter, preservation of updates implies preservation of reads.
Indeed, for any state $c$, applying the first law before the change
(\eqref{eq:1.16}), preservation of updates, and the second law after the change gives
\[
\begin{aligned}
g'(h(c))
&=g'\bigl(h(p(c,g(c)))\bigr)\\
&=g'\bigl(p'(h(c),u(g(c)))\bigr)\\
&=u(g(c)).
\end{aligned}
\]
This derivation does not require $h,u$ to be invertible.
The following construction expresses preservation of reads and updates by the same
maps as one commutative diagram.

\begin{proposition}[A square combining the two operations]\label{prop:4.37}

Combine reads and updates into a single map of sum types,
\begin{equation*}
c_L:C\sqcup(C\times V)\longrightarrow V\sqcup C,\qquad
c_L|_C=g,\quad c_L|_{C\times V}=p
\tag{4.45}\label{eq:4.45}
\end{equation*}

where the right-hand side uses the inclusion into the corresponding summand.
The following equation is equivalent to \eqref{eq:4.44}:
\begin{equation*}
(u\sqcup h)c_L
=c_{L'}\bigl(h\sqcup(h\times u)\bigr).
\tag{4.46}\label{eq:4.46}
\end{equation*}

\end{proposition}

\begin{proof}

Restriction to the left input summand $C$ gives the read equation, and restriction
to the right input summand $C\times V$ gives the update equation.
Conversely, equality on both summands gives equality of maps on the sum.

\end{proof}

In \eqref{eq:4.46}, the same $h$ is used for the input of reads, the state
component of the input of updates, and the output of updates.
This is the preservation condition for the shared state.

\begin{proposition}[Invertible changes preserving both operations]\label{prop:4.38}

The group of invertible changes of a lens is the following subgroup of
$\mathrm{Aut}(C)\times\mathrm{Aut}(V)$:
\begin{equation*}
\begin{aligned}
G_L
 &=G_{\mathrm{get}}\cap G_{\mathrm{put}}
 =G_{\mathrm{put}}\subseteq G_{\mathrm{get}},\\
G_{\mathrm{get}}&=\{(h,u)\mid gh=ug\},\\
G_{\mathrm{put}}&=\{(h,u)\mid hp=p(h\times u)\}
\end{aligned}
\tag{4.47}\label{eq:4.47}
\end{equation*}

Equation~\eqref{eq:4.46} says that the input and output changes constructed from
this same pair $(h,u)$ commute with $c_L$.

\end{proposition}

\begin{proof}

Both conditions are closed under identities, composition, and inverses.
Their intersection consists exactly of the invertible morphisms satisfying
\eqref{eq:4.44}. Since preservation of updates implies preservation of reads,
this intersection equals $G_{\mathrm{put}}$.
The last assertion follows from Proposition~\ref{prop:4.37}.

\end{proof}

In Example~\ref{ex:1.35} of Chapter~\ref{chap:1}, of the four bijections
preserving reads, only two also preserved updates.
Thus the inclusion in \eqref{eq:4.47} is strict in general.
Preserving updates further constrains the allowed correspondence of states.

\begin{example}[A semantics-preserving morphism that merges information]\label{ex:4.39}

For product lenses $C=V\times K$ and $C'=V\times K'$ over a nonempty view set $V$,
any map $t:K\to K'$ satisfies \eqref{eq:4.44} through
\begin{equation*}
h(v,k)=(v,t(k)),\qquad u=\mathrm{id}_V
\tag{4.48}\label{eq:4.48}
\end{equation*}

Indeed, both reads give $v$, and an update to the displayed value $w$ gives
$(w,t(k))$ in either order.
If $K=\{0,1\}$ and $K'=\{*\}$, then $h$ is not injective.
General semantics-preserving morphisms include such aggregation of information;
taking the invertible ones yields the group in \eqref{eq:4.47}.

\end{example}

\section{Carrying protocol operations and adapters}\label{sec:4.11}

In a protocol, what is shared are the state sets at the vertices.
The same state map is used for the operations that enter a vertex and for those that
leave it.
Furthermore, if an adapter connects two protocols, that connection is preserved at the
same time.

\begin{proposition}[Preservation from named operations to all executions]\label{prop:4.40}

Take realizations $X,X'$ of $\mathrm{Prot}(Q,\Pi,O)$ of Chapter~\ref{chap:1} and vertex
maps $a_v:X(v)\to X'(v)$.
If, for each named edge $e:v\to w$ and each observation,
\begin{equation*}
a_wX(e)=X'(e)a_v,\qquad
o_{X'}(v)a_v=o_X(v)
\tag{4.49}\label{eq:4.49}
\end{equation*}
then $a$ determines a unique semantics-preserving morphism of protocols.
Moreover, for adapters $q:X\to Y$ and $q':X'\to Y'$ and semantics-preserving morphisms
$a:X\to X'$ and $b:Y\to Y'$,
\begin{equation*}
bq=q'a
\quad\Longleftrightarrow\quad
b_vq_v=q'_va_v\quad\text{(for every vertex }v\text{)}
\tag{4.50}\label{eq:4.50}
\end{equation*}
\end{proposition}

\begin{proof}

Condition~\eqref{eq:4.49} is the condition of Proposition~\ref{prop:1.38} and gives a
unique semantics-preserving morphism.
Reading the preservation condition for adapters in \eqref{eq:1.23} componentwise as a
natural transformation gives \eqref{eq:4.50}.

\end{proof}

In this check, operations with different names are treated as separate generating edges
even when their actions coincide.
When adapters are changed in succession, pasting the two commutative squares shows that
\eqref{eq:4.50} also holds for the composed change.

\section{Comparison of semantics and typed constructions}\label{sec:4.12}

Invertible changes that preserve a semantics-preserving morphism of lenses or a
protocol adapter can be recovered after passing to the typed constructions.
Common to both is the following correspondence through a fully faithful functor.

\begin{theorem}[Preserving comparisons between semantics and typed constructions]\label{thm:4.41}

For a fully faithful functor $F:\mathcal{C}\to \mathcal{D}$ and a morphism $c:X\to Y$,
define the group preserving the comparison by
\begin{equation*}
\mathrm{Cmp}(c)=
\{(a,b)\in\mathrm{Aut}_{\mathcal{C}}(X)\times\mathrm{Aut}_{\mathcal{C}}(Y)
\mid bc=ca\}
\tag{4.51}\label{eq:4.51}
\end{equation*}
Then
\begin{equation*}
\mathrm{Cmp}(c)\xrightarrow{\sim}\mathrm{Cmp}(F(c)),
\qquad(a,b)\longmapsto(F(a),F(b))
\tag{4.52}\label{eq:4.52}
\end{equation*}
is a group isomorphism.
If $c$ is an isomorphism, the inverse of the first projection,
$a\mapsto(a,cac^{-1})$, also commutes with this correspondence.

\end{theorem}

\begin{proof}

A functor preserves the equation $bc=ca$, so the map is defined and is a group
homomorphism.
By fullness, the automorphisms of $F(X),F(Y)$ and their inverses lift to the original
category.
By faithfulness, the composites of the lifted morphisms are identities, and
$bc=ca$ also follows from $F(b)F(c)=F(c)F(a)$.
Hence the map is surjective. Injectivity also follows from faithfulness.
When $c$ is an isomorphism, $b=cac^{-1}$, and the last claim follows because a functor
preserves composites and inverses.

\end{proof}

Proposition~\ref{prop:1.43} of Chapter~\ref{chap:1} gives fully faithful functors from
the semantics of fixed-view lenses and of fixed protocols to the categories of typed
objects constructed from each.
Therefore the invertible changes that preserve semantics-preserving morphisms of
lenses or protocol adapters can be recovered exactly from the typed constructions.
For lenses whose views also change, take the view component in
Definition~\ref{def:1.42} to be $u$ and the write component to be $h\times u$.
By \eqref{eq:4.44} and Proposition~\ref{prop:4.37}, extracting and assembling
the role maps are mutually inverse, giving the same conclusion.

Once a category with these semantics-preservation conditions and a functor to the
base have been defined, choose strongly opcartesian edge lifts and specified comparisons.
The defects of \S\S\ref{sec:4.5}--\ref{sec:4.7} then express discrepancies
between changes that preserve the actual operations.
To connect further with transport of readings equipped with geometry, supply
the core morphism and the comparisons of local realizations of Definition~\ref{def:4.9}.
Providing these data in this order connects the meaning of operations, changes of
reading, and lifts of geometries.

\section*{Summary of the Chapter}

In this chapter we constructed transport along changes of reading, and gave conditions
under which several transports and specified comparisons become coherent for the same
selection of edges.

\begin{itemize}
\item \textbf{Existence and universality of transport.} We constructed the core from an
exact change of extraction and showed that its canonical morphism is strongly
opcartesian.
We also constructed the canonical transport of geometries and, for general morphisms of
cores, extracted the three comparisons of local realizations as the existence condition
(\S\S\ref{sec:4.2}--\ref{sec:4.3}).
\item \textbf{Coherence of units and composition.} From the uniqueness of factorizations
we obtained the functors between fibers and the natural isomorphisms for units and
composition. Triple composites and the projection to the core
are coherent as well (\S\ref{sec:4.4}).
\item \textbf{Obstructions for specified comparisons.} We expressed the difference from
the canonical comparison as a noncommutative automorphism and showed that vanishing of the obstruction is equivalent to an edge
reselection making all faces coherent.
Pasting faces requires additional coherence conditions among the specified comparisons
(\S\S\ref{sec:4.5}--\ref{sec:4.8}).
\item \textbf{Simultaneous vanishing across levels.} After the selection at the core has
been lifted to the geometry, being able to remove the remainder by kernel changes for the
same selection corresponds to overall coherence.
Even if each level has a solution independently, without a lift there is no overall
solution (\S\ref{sec:4.9}).
\item \textbf{Comparisons preserving the meaning of operations.} For lenses,
reads and updates use the same state map, and preservation of updates also implies
preservation of reads. For protocols, all named operations and adapters are preserved
simultaneously. The invertible changes preserving these comparisons can be recovered
across the constructed typed categories (\S\S\ref{sec:4.10}--\ref{sec:4.12}).
\end{itemize}

For the order-processing API of the opening, the discrepancy between the route that moves
directly from the old version to the new one and the route through the intermediate
version is extracted as a defect.
Even a discrepancy that could be absorbed into one of two independent edges must be made
coherent simultaneously once that edge is shared with another face.
Theorem~\ref{thm:4.23} characterizes the existence of a family of transformations
realizing the specified comparisons on every face under consideration as the vanishing
of the obstruction.
When both levels, core and geometry, are treated, the lifts and the kernel changes are
again performed for the same selection.

In the next chapter we compare the route that changes the base of the input before
constructing and the route that constructs first and then changes.
